\section{CONCLUSION}
\label{sec:conclusion}
We investigated the motion of a ferromagnetic solid sphere rolling without slipping on an incline in an external magnetic field. Our simulations show that the downhill trajectory of such a sphere is considerably influenced by a magnetic field as small as that of the Earth. In general, a straight downhill trajectory acquires a transverse component similar to that of the Lorentz force. However, the deflection is highly dependent on the initial direction and magnitude of both $\bm{\mu}_0$ and $\bm{B}$. 

Surprisingly, we did not observe any symmetry on exchanging the data of $\bm{B}$ and $\bm{\mu}_0$ when investigating the influence of the azimuthal angular deviation from alignment of $\bm{B}$ and $\bm{\mu}$ despite the fact that exchanging them would only change the sign of the magnetic torque. However, a mirror symmetry emerges in the $y(t)$-coordinate when the polar angle is varied under this exchange. Moreover, we could identify three qualitatively different trajectories: Curved downward trajectories not twisting around the $x=0$ line for low values of $\mu_0$ and $B$, trajectories twisting around the $x=0$ line for high values of $B$ and high azimuthal deviations of $\bm{B}$ from $\bm{\mu}_0$, and trajectories looping around or oscillating at the origin for high values of $\mu_0$ and high values of azimuthal or polar deviation. 

To conclude, the apparent simplicity of the system paired with an intricate dependence on initial values turns it into an excellent proof-of-concept for the Lagrangian-d'Alembert approach to non-holonomic, dissipative systems. We demonstrated that even with multiple dissipative forces, the equations of motions can be solved numerically with sufficient stability, and provided a software to that end. In the future, it may be fruitful to investigate the system for chaotic behavior. Since the approach is also valid for $\varphi = 0^\circ$, i.e., a flat surface, comparison with analytical solutions available in certain such cases may be of interest.